\documentclass[11pt,a4paper]{article}
\usepackage[T1]{fontenc}
\usepackage[utf8]{inputenc}
\usepackage[french]{babel}
\frenchsetup{SmallCapsFigTabCaptions=false}
\usepackage{lmodern,microtype,amsmath}
\usepackage[margin=22mm]{geometry}
\usepackage{booktabs,array,tabularx,longtable,enumitem,xcolor,tikz}
\usetikzlibrary{arrows.meta,positioning}
\usepackage{hyperref}
\hypersetup{colorlinks=true,linkcolor=black,citecolor=black,urlcolor=black,
 pdftitle={Méthodes sans LLM pour LANCE : références fortes et programme de mesure},
 pdfauthor={Codex, synthèse de trois recherches et contre-relecture indépendante},pdflang={fr-FR}}
\urlstyle{same}
\def\UrlBreaks{\do\/\do\-\do\_\do\.\do\:}
\newcolumntype{Y}{>{\raggedright\arraybackslash}X}
\newcommand{\code}[1]{\nolinkurl{#1}}
\newcommand{\revision}{2026-10-07 14:44 (Europe/Brussels, UTC+02:00)}
\setlength{\parskip}{3pt}
\setlength{\emergencystretch}{3em}
\setlist{nosep,leftmargin=5mm,topsep=3pt}
\clubpenalty=10000 \widowpenalty=10000
\setcounter{tocdepth}{1}
\definecolor{line}{HTML}{C9CED5}
\definecolor{ink}{HTML}{252A32}
\tikzset{box/.style={draw=line,rounded corners=2pt,align=center,text=ink,
 text width=42mm,minimum height=15mm,inner sep=5pt,font=\small},
 flow/.style={-{Latex[length=2mm]},draw=ink,line width=.6pt}}
\title{Méthodes sans LLM pour LANCE\\
 \large Références fortes et programme de mesure pour le comité de thèse}
\author{Synthèse Codex de trois recherches spécialisées\\
 \small Projet LANCE ; contre-relecture par des agents distincts des auteurs}
\date{Dernière révision : \revision\\[3pt]
 \small Recherche exploratoire et protocole proposé ; aucun gain expérimental établi}
\begin{document}
\maketitle
\begin{abstract}
La réussite d'un audit utilisant un modèle de langage ne mesure pas l'apport
du modèle. Cette recherche examine les références sans LLM qui permettraient
de répondre à l'objection du comité de thèse. Trois revues ciblées étudient les
méthodes déterministes, le raisonnement classique et les expériences recevables.
La priorité proposée est une référence conditionnelle compétente, évaluée sur
des dossiers et critères communs ; une procédure fixe reste un contraste
secondaire. Jointures, moteurs de politiques et vérification de configurations
sont retenus seulement si les données et la propriété les justifient. Six
expériences hors ligne séparent qualité des conclusions, fausses alertes,
contexte, robustesse et utilité du rapport. Une contre-relecture impose un gel
de la référence, un univers indépendant des candidats LANCE et des coûts humains
complets. Les documents décrivent une méthode de mesure, sans résultat de
supériorité ni d'équivalence ; une décision favorable aux règles seules reste
scientifiquement recevable.
\end{abstract}
\tableofcontents
\clearpage

\section{Question : mesurer un effet marginal utile}
LANCE réalise des audits de sécurité d'objets connectés (IoT). Il combine des
outils, des règles et des modèles de langage (LLM). La question du comité est :
\textbf{sur quelles tâches le LLM apporte-t-il un bénéfice utile,
au-delà d'une méthode simple correctement conçue, et à quel coût ?} Déclarer
que le système complet trouve des problèmes ne répond pas à cette question.

La contribution de thèse envisageable devient une décomposition empirique :
quelles décisions sont déjà résolues par les règles, où subsiste une difficulté,
et si le modèle améliore cette difficulté avec des preuves vérifiables. Cette
formulation est une proposition de recherche, pas une revendication de nouveauté
mondiale. Le protocole historique distingue déjà interprétation, adaptation et
restitution \cite{historique} ; l'approfondissement précise surtout la force
des comparateurs classiques et les conditions de leur sélection. SQL désigne
le langage de requêtes relationnelles ; Datalog, un langage de règles logiques.
OPA est l'Open Policy Agent, et Cedar un langage de politiques. HSA signifie
\emph{Header Space Analysis}. Les autres noms du tableau sont des outils ou
des formats, décrits dans les rapports spécialisés.

\textbf{Sans LLM ne signifie pas sans adaptation.} Un programme peut utiliser
des conditions, de l'état, des relations et une politique d'abstention sans
apprentissage. Sans LLM ne signifie pas non plus sans IA : l'apprentissage
statistique classique appartient à cette catégorie, mais demande des données
et une validation qui ne le rendent pas automatiquement plus simple.

\section{Méthode et statut des preuves}
\subsection{Recherche organisée et lectures}
L'étude conserve le dossier commencé le 27 septembre 2026. L'approfondissement
du 7 octobre répartit trois tâches entre agents : outils déterministes
\cite{outils}, raisonnement classique \cite{raisonnement}, et programme
d'expériences \cite{experiences}. Des agents distincts des auteurs assurent la
contre-relecture du texte et des pages rendues. Le journal \code{validation.md}
identifie les versions, objections, corrections et limites réellement relues.

Les sources sont primaires : spécifications, documentations officielles et
articles des auteurs. La sélection est ciblée, non systématique ; elle vise les
entrées nécessaires, la transformation produite, la portée des conclusions et
les conditions d'évaluation. Chaque registre spécialisé précise les sections
lues, date ou version, consultation et limites. Les articles ne sont pas
reproduits ; les résultats de corpus hétérogènes ne sont pas classés ensemble.
Le registre \code{sans-llm-sources.json} réunit ces lectures et les sources
complémentaires de la présente synthèse.

La lecture NIST SP~800-115, sections~3.3--3.4, justifie notamment une référence
de revue de règles et de configurations, avec collecte et accès explicités
\cite{nist115}. SP~800-53A, section~2.4, relie objets, méthode, couverture et
critère d'évaluation \cite{nist53a}. Leur transfert aux fiches LANCE est notre
proposition ; il ne constitue ni certification, ni validation de LANCE.

\subsection{État local et faits observables}
Le commit de base est \code{a8dbeb9404f0cad960d61e42f02a1d70ca143a59}, avec
modifications et conflits locaux. Les empreintes des lectures figurent dans
les registres spécialisés. La lecture du code est statique ; aucun import
applicatif, campagne, appel LLM de benchmark ou collecte réseau n'a été réalisé.
Les conflits empêchent de présumer que cette copie s'exécute ou représente le
déploiement. Les résultats historiques restent attribués à leurs dates.

Nous distinguons \emph{fait sourcé}, \emph{fait de lecture statique},
\emph{hypothèse} et \emph{protocole proposé}. La compilation et l'inspection
des PDF vérifient les documents ; elles ne mesurent pas la valeur du modèle.

\section{Références simples : sélectionner la bonne tâche}
\subsection{Ne pas confondre les objets évalués}
Un inventaire de service, un avis concernant un produit, une règle de
configuration et une observation de comportement établissent des faits
différents \cite{outils}. Une correspondance de version peut justifier une
applicabilité selon un avis ; elle ne prouve pas un impact observé. Une
politique d'autorisation calcule un droit attendu ; sa décision ne prouve pas
que l'équipement applique ce droit. Une vérification formelle est exacte sur
son modèle, sous réserve que ce modèle décrive correctement les données.

Le critère d'audit doit donc préciser propriété, ressource, identité, contexte,
période et preuve requise. L'absence d'une donnée ne devient pas un état sain.
L'absence d'un tuple SQL, une valeur \code{NULL} ou une règle non déclenchée
n'établit pas à elle seule la complétude. Une déclaration explicite de couverture
et une politique d'abstention sont nécessaires \cite{raisonnement}.

\subsection{Ordre proposé et conditions d'entrée}
Le tableau~\ref{tab:choix} propose une progression. L'ordre exprime la
compatibilité avec les dossiers disponibles et la lisibilité des erreurs,
sans estimation chiffrée de charge ni classement de performances.

\begin{table}[htbp]\centering\small
\begin{tabularx}{\linewidth}{p{32mm}YY}
\toprule
Référence & Ce qu'elle permet de tester & Condition / décision proposée\\
\midrule
Checklist fixe, D0 & Application systématique de critères explicites. & Contraste secondaire ; critères compétents, inconnus conservés.\\
Règles conditionnelles & Contexte, exceptions, contradictions et abstention. & Référence principale ; D1 existant ou lecteur hors ligne proposé, selon tâche.\\
Jointures SQL & Attribution, rapprochement de pièces et traçabilité. & Ajouter si les relations l'exigent ; même politique de décision.\\
Datalog / Soufflé & Interdépendance des règles et provenance d'inférences. & Option après SQL ; la récursion existe aussi en SQL.\\
Table de droits, OPA ou Cedar & Calcul d'une autorisation attendue et comparaison à une observation. & Politique propriétaire disponible ; table simple avant moteur spécialisé.\\
Matching inventaire--avis & Applicabilité sous identité et conditions connues. & Inventaire crédible et avis figés ; abstention si contexte insuffisant.\\
OVAL / OpenSCAP & Conformité à un profil sur export local compatible. & Profil adapté et mêmes pièces locales dans les deux bras.\\
FIREMAN, HSA, Batfish & Propriétés d'un modèle de configuration réseau. & Piste séparée ; exports supportés et question précise indispensables.\\
\bottomrule
\end{tabularx}
\caption{Références proposées par tâche. Les outils externes et le lecteur
hors ligne ne sont pas annoncés comme intégrés dans LANCE.}\label{tab:choix}
\end{table}

Nmap ou des matchers et workflows de contrôles bénins constituent des
implémentations classiques pertinentes pour certaines propriétés
\cite{outils}. Comparer leur utilité complète à LANCE peut être utile, à
condition de publier leurs couvertures et accès respectifs. Cela ne suffit
pas à isoler l'effet du LLM. OpenSCAP et les avis liés aux inventaires logiciels
SBOM (\emph{Software Bill of Materials}) et au format CSAF
(\emph{Common Security Advisory Framework}) changent aussi
l'information disponible lorsque les configurations ou inventaires logiciels
ne sont pas remis aux deux bras.

Le premier pilote écarte l'apprentissage classique faute de labels
indépendants identifiés, les moteurs spécialisés firmware faute d'artefacts
adaptés, et NetPlumber faute de question sur les mises à jour. Ces exclusions
sont des décisions de périmètre ; elles ne prouvent aucune faiblesse de ces
méthodes. Leur retour devient recevable lorsqu'une question et ses données
sont disponibles \cite{raisonnement}.

\section{D1 existe déjà : préserver les contrastes}
\textbf{Fait de lecture statique.} D1 désigne la politique d'automatisation
conditionnelle sans LLM ; A1, la politique adaptative avec LLM. Le contrat
\code{analysis-verification} fixe un inventaire public, le profil
\code{full} et les phases~3--4 : analyse et vérification. Le scanner initial
est partagé ; les règles participent donc déjà au succès d'A1. D1 n'est pas
une référence couvrant toutes les phases de l'audit. Les sources visibles
et limites de ce diagnostic sont détaillées dans \cite{experiences}.

Les identifiants D0, A0 et D1-R du protocole historique restent des variantes
proposées, à revalider avant emploi. Un \textbf{lecteur à règles sur dossiers}
est ici un adaptateur hors ligne proposé ; il ne doit pas être présenté comme
le D1 actuellement implémenté. De même, une fonction de rapprochement d'avis
de vulnérabilité CVE (\emph{Common Vulnerabilities and Exposures}) visible
sur un chemin \code{compact} ne prouve pas sa couverture dans l'expérience
\code{full}. Ce constat est limité au chemin inspecté.

Les connaissances publiques spécifiques aux simulateurs demandent un audit
de fuite. Le partage d'un oracle de production peut aussi produire des erreurs
communes : le contrat de preuve LANCE ne devient pas la vérité terrain du test.
Les identifiants et empreintes déclarés aident à comparer les artefacts ; ils
ne certifient pas l'identité de l'état réellement observé.

\subsection{Deux effets à garder séparés}
Le premier contraste mesure \textbf{l'interprétation de pièces communes} :
normalisation et contexte accessibles identiques, critère identique, lecteur à
règles contre lecteur LLM. Il ne mesure pas la collecte. Le second mesure
\textbf{l'utilité d'un système complet}, avec ses outils, accès et couverture
déclarés. Il peut guider un choix pratique, mais l'attribution au LLM y demeure
plus faible. Le contraste D1/A1 phases~3--4 est un ajout du modèle au socle
commun ; des observations finales différentes empêchent de le réduire à une
comparaison d'interprètes.

Même sans recherche externe, le préentraînement du LLM constitue une
connaissance supplémentaire. Des noms et avis fictifs cohérents permettent un
contrôle de mémorisation, sans prouver l'absence de toutes les fuites. Une date
récente ne garantit pas que l'information soit absente du modèle. Un résultat
doit donc nommer l'effet du composant LLM, plutôt que revendiquer automatiquement
un raisonnement pur.

La figure~\ref{fig:attribution} montre le contraste hors ligne recommandé.
La sélection des pièces est indépendante des candidats produits par LANCE,
afin de conserver les omissions dans l'univers évalué.
\begin{figure}[htbp]\centering
\begin{tikzpicture}
\node[box] (data) {Dossiers et contexte\\figés, communs};
\node[box,right=9mm of data,yshift=13mm] (rules) {Lecteur à règles\\accepté sur développement};
\node[box,right=9mm of data,yshift=-13mm] (llm) {Lecteur LLM\\configuration figée};
\node[box,right=9mm of rules,yshift=-13mm] (oracle) {Oracle indépendant\\et coûts complets};
\draw[flow] (data.east) -- (rules.west);
\draw[flow] (data.east) -- (llm.west);
\draw[flow] (rules.east) -- (oracle.west);
\draw[flow] (llm.east) -- (oracle.west);
\end{tikzpicture}
\caption{Protocole proposé pour attribuer une différence d'interprétation
au composant LLM. Le modèle peut conserver des connaissances préentraînées.
La figure est modifiable dans cette source.}\label{fig:attribution}
\end{figure}

\section{Six expériences et premières priorités}
Les fiches complètes \cite{experiences} décrivent données communes, oracle,
contrôles, comparateurs, métriques, coûts et critères d'arrêt. Les expériences
ci-dessous sont \textbf{proposées}, sans corpus admissible ni effectif validé.

\begin{longtable}{>{\raggedright\arraybackslash}p{25mm}>{\raggedright\arraybackslash}p{65mm}>{\raggedright\arraybackslash}p{52mm}}
\caption{Programme de tests hors ligne ; priorités décidées avant lecture du test réservé.}\label{tab:tests}\\
\toprule
Expérience & Contraste et exemple de dossier & Décision ou limite visée\\
\midrule\endfirsthead
\toprule Expérience & Contraste et exemple de dossier & Décision ou limite visée\\\midrule\endhead
Routine & Critère direct, pièce suffisante ou absente ; procédure fixe, règles conditionnelles et lecteur LLM. & Les règles suffisent-elles ? Ne pas rendre la procédure fixe artificiellement faible.\\
Cas sains / faible prévalence & Positifs, négatifs proches et propriété non documentée ; les cas sains ressemblent aux alertes. & Fausses alertes, abstention et couverture ; \textbf{priorité avec routine}.\\
Autorisation & Même réponse avec identité et politique différentes ; politique absente ou contradictoire. & Table de droits avant OPA/Cedar ; gain utile seulement si contexte correctement relié.\\
Information répartie & Pièces séparées, dates incompatibles ou contradictions ; jointures simples contre lecteur LLM. & Mesurer liaison et preuve, pas seulement accumulation de texte ; prioritaire si pièces disponibles.\\
Variations équivalentes & Même sémantique, ordre des champs ou formulation modifiés ; contenu incomplet séparé. & Distinguer robustesse du parseur et interprétation ; pas de handicap sélectif pour les règles.\\
Restitution humaine & Même ensemble de faits validés, rapport déterministe contre rédaction LLM ; revue de décisions par lecteurs. & Temps de vérification, erreurs et utilité ; gain rédactionnel distinct de la détection.\\
\bottomrule
\end{longtable}

\textbf{Commencer par routine et cas sains.} Ces tests donnent une première
estimation de la qualité, des inconnus et du coût, avec les références simples
les plus crédibles. Ajouter ensuite autorisation ou information répartie
si des pièces suffisantes existent. La restitution demande une procédure
humaine séparée : ordre contrebalancé, décisions et temps mesurés, exposition
répétée au même cas contrôlée. Aucun effectif ni seuil n'est fixé arbitrairement.
Geler les faits sources ; conserver et compter les omissions et inventions
dans les rapports générés. Une variation de format sur le même corpus peut
compléter le pilote : l'échec de parsing reste une sortie mesurée dans la
strate brute, même s'il invalide le contraste sur dossiers normalisés.

Pour un rejeu conditionnel, toutes les branches enregistrées ne suffisent
pas lorsque les réponses dépendent de l'historique. Restreindre le pilote à
des lectures sans effet et indépendantes de cet historique, ou à des histoires
fournies cohérentes avec leur état ; sinon la tâche reste hors couverture.
Les exemples synthétiques servent aux contrôles ciblés et sont signalés.
Ils ne deviennent pas des résultats terrain.

\section{Ce qu'il faut mesurer pour pouvoir conclure}
\subsection{Oracle, univers et sorties}
Définir l'univers indépendant de propriétés et dossiers \emph{avant} les
candidats LANCE : inclure ce que son collecteur a omis. Sur un corpus limité
aux candidats, le rappel est conditionnel et ne décrit pas l'audit complet.
Pour chaque cas, l'oracle indépendant décrit deux objets : $T_{\text{état}}$,
état de la propriété, et $T_{\text{preuve}}$, suffisance des preuves
accessibles. Une propriété réellement violée peut
rester non démontrable dans le dossier ; la bonne sortie est alors l'incertitude,
selon le contrat préenregistré. Le problème d'oracle reste une limite générale
des tests logiciels \cite{experiences}.

Prévoir une revue distincte, des justifications, le traitement des désaccords
et le masquage de l'identité du système lorsque possible. Les évaluateurs ne
réutilisent pas seulement le verdict automatique de LANCE. Conserver
contradictions, erreurs techniques, dépassements et essais incomplets : les
retirer peut transformer la fragilité en gain apparent.

Les mesures comprennent précision des conclusions étayées, rappel dans
l'univers annoncé, fausses alertes sur cas sains, couverture des décisions,
abstention justifiée et taux d'échec. Rapporter aussi la qualité des références
aux pièces et le nombre de conclusions sans support. Une seule moyenne masque
des erreurs différentes ; stratifier par famille préenregistrée et afficher les
dénominateurs. Le rappel sur les états réellement positifs de l'univers
indépendant et le rappel conditionnel aux positifs démontrables sont publiés
séparément. Les cas non démontrables, abstentions et échecs restent dans les
totaux et la couverture. Des strates trop petites restent descriptives.
Dans une strate saine, privilégier le taux de fausses alertes avec nombre
de cas réellement sains au dénominateur ; le score F1 peut être nul ou
indéfini selon les sorties et la convention, sans décrire seul cette qualité.

\subsection{Indépendance et incertitude}
Comparer les bras sur mêmes familles et mêmes conditions. Des répétitions du
modèle sur une famille ne deviennent pas autant de cas indépendants. Une
famille d'origine regroupe les cas partageant générateur, implémentation,
service ou structure ; leurs variantes restent dans le même groupe. Les six
expériences du tableau ne sont pas six unités indépendantes. Le test
réservé doit différer des données utilisées pour écrire critères, règles,
prompts et sélection de modèle. Le découpage public S1--S19 / S20--S29 ne
garantit pas à lui seul l'indépendance historique de ces familles.
Un test réservé aux auteurs ne certifie pas son absence du préentraînement ;
de nouvelles familles privées réduiraient ce risque sans le faire disparaître.

Le pilote est d'abord descriptif. Le choix des marges utiles, de la mesure
principale, des garde-fous et de l'effectif précède une étude confirmatoire et
dépend de la variabilité observée et de la structure des familles. Une différence
non significative ne prouve pas une équivalence \cite{lakens}. Une conclusion
d'équivalence ou de non-infériorité demande un protocole approprié et des bornes
justifiées. Un intervalle large appelle une conclusion indéterminée. Il n'est
pas proposé d'appliquer mécaniquement un test prévu pour moyennes indépendantes
à tous les scores ou données groupées.

\section{Coûts et décisions recevables}
\subsection{Mesurer le travail déplacé}
Les règles coûtent en critères, normalisation et maintenance ; le modèle coûte
en intégration, appels, validation et correction. Les données, oracles et
profils nécessaires doivent être comptés pour chaque bras. Mesurer séparément
temps de développement, entretien par changement, temps de revue humaine,
calcul, latence et consommation. Ne pas additionner directement heures, euros
et tokens. Une monétisation éventuelle précise taux, prix et date ; l'absence
de tokens ne signifie pas absence de travail.

Sur un horizon annoncé de $n$ dossiers, une représentation proposée est
\[
C(n)=C_{\text{développement}}+C_{\text{maintenance}}(n)
       +\sum_{i=1}^{n}(C_{\text{exécution},i}+C_{\text{revue},i}).
\]
Tous les termes doivent employer la même unité si un total est calculé.
Publier aussi les mesures brutes et les coûts partagés : un meilleur total
peut dépendre de l'horizon et du prix retenus. Aucun point de rentabilité
n'est déduit sans mesure.

\subsection{Décider sans obliger le LLM à gagner}
Le tableau~\ref{tab:decision} donne les issues admissibles. Ces conditions
doivent être recevables avant de connaître le résultat du test.
\begin{table}[htbp]\centering\small
\begin{tabularx}{\linewidth}{YY}
\toprule
Résultat possible, à établir & Conséquence proposée\\
\midrule
Règles adéquates, coût inférieur, preuve de qualité suffisante & Conserver la méthode classique pour la tâche.\\
Gain LLM utile, garde-fous respectés, coût acceptable & Retenir le modèle pour cette tâche et les conditions mesurées.\\
Gain seulement sur rédaction et revue & Limiter la conclusion à la restitution ; garder les faits validés comme entrée commune.\\
Gain disparaissant avec règles compétentes ou avis fictifs & Réviser l'attribution ; coût de la connaissance ou mémorisation à distinguer.\\
Incertitude trop large, oracle ou corpus inadmissible & Ne conclure ni gain ni équivalence ; corriger la condition manquante.\\
\bottomrule
\end{tabularx}
\caption{Règles de décision proposées. « Adéquat » et « utile » demandent
des seuils justifiés et fixés avant le test confirmatoire.}
\label{tab:decision}
\end{table}

Un fonctionnement mixte par famille peut être envisagé seulement si son
routage est fixé sur développement et évalué ensuite. Choisir pour chaque
cas le vainqueur après lecture du test ne constitue pas un système déployable.
La référence principale est sélectionnée et gelée avant test parmi les méthodes
admissibles, sans sélection opportuniste a posteriori.

\section{Travail à préparer et limites}
\subsection{Lots concrets avant une expérience}
Le premier lot est un \textbf{registre d'éligibilité} : propriété, pièces,
provenance, accès, contexte, indépendance, oracle et inconnus pour chaque famille.
Il permet de refuser une comparaison non identifiable. Le deuxième lot est
une \textbf{référence acceptée sur développement} : critères revus, contextes,
négatifs proches, inconnus, justifications, effort humain enregistré. Geler
code, corpus, avis, prompts, modèles, normalisation et configuration avant
de découvrir le test ; toute nouvelle version nécessite un nouveau test réservé.

Le troisième lot est un \textbf{pilote apparié hors ligne}, avec routine et
cas sains, puis contextualisation si admissible. Sa sortie est un diagnostic
des erreurs, coûts et variabilité, pas une revendication confirmatoire. Le
quatrième lot prépare le \textbf{test réservé} : marges utiles, plan d'analyse,
familles réellement indépendantes, garde-fous et effectif justifié. La revue
humaine de restitution est un lot séparé si la question devient prioritaire.
Ces lots restent proposés ; aucune intégration n'est annoncée achevée.

\subsection{Objections de la contre-relecture}
La première contre-revue a exigé six corrections : distinguer D1 implémenté
du lecteur proposé ; fixer une porte d'acceptation de la référence ; qualifier
la connaissance préentraînée ; garantir un rejeu cohérent ; inclure les omissions
du collecteur ; laisser la pagination suivre le contenu. Le journal conserve
également les corrections de la contre-revue expérimentale : labels d'état et
de preuve séparés, matrice normative légitime, erreurs de restitution conservées,
conventions de F1 et échecs du parseur explicités. Il identifie
les arbitrages et les versions effectivement inspectées. Les conditions non
résolues concernent encore les données éligibles, l'oracle indépendant,
l'indépendance des familles et les seuils utiles. Aucune relecture documentaire
ne peut les déclarer expérimentalement satisfaites.

\subsection{Limites de la conclusion}
La recherche est ciblée et centrée sur des dossiers défensifs ; elle ne décrit
pas toutes les méthodes d'audit IoT et ne mesure pas la découverte autonome.
Les outils externes restent non exécutés, leurs compatibilités non établies et
les coûts qualitatifs. Un gain sur dossiers synthétiques ne prouve pas un gain
sur réseaux réels ; un gain d'interprétation ne prouve pas une meilleure collecte.
La copie conflictuelle limite aussi les conclusions sur l'état exécutable.

\section{Conclusion}
La première réponse au comité est de construire une mesure qui pourrait
conclure en faveur des méthodes simples. Le meilleur point de départ est une
référence conditionnelle compétente, des dossiers communs et un oracle distinct
du verdict produit par LANCE. Routine et cas sains précèdent les formalismes
plus complexes ; jointures et politiques interviennent quand les pièces les
justifient. La plus-value du LLM reste une hypothèse à tester. Démontrer qu'il
est inutile sur certaines tâches, ou utile seulement pour la restitution,
serait également une conclusion exploitable pour la thèse.

\phantomsection
\addcontentsline{toc}{section}{Références}
\begin{thebibliography}{9}\small
\bibitem{historique} Projet LANCE, \emph{Apport de l'agent face à l'automatisation},
étude commencée le 27 septembre 2026, rapport révisé le 1er octobre 2026 à 10:23.
\href{report.pdf}{Rapport historique}, \href{protocol.md}{protocole} et
\href{state-of-the-art.md}{état de l'art ciblé}. Lecture ciblée des comparateurs
et conditions ; résultats historiques non reproduits.
\bibitem{outils} Agent \code{sans_llm_outils}, \emph{Méthodes simples sans LLM
pour l'audit IoT : références défensives et comparateurs pour LANCE},
7 octobre 2026. \href{sans-llm-outils.pdf}{Rapport} et
\href{sans-llm-outils-sources.json}{registre des lectures}.
\bibitem{raisonnement} Agent \code{sans_llm_raisonnement}, \emph{Raisonnement
sans LLM pour l'audit IoT : références simples, politiques et vérification de
configurations}, 7 octobre 2026. \href{sans-llm-raisonnement.pdf}{Rapport} et
\href{sans-llm-raisonnement-sources.json}{registre des lectures}.
\bibitem{experiences} Agent \code{sans_llm_experiences}, \emph{Mesurer la valeur
marginale des LLM face à des audits simples : diagnostic de D1 et six expériences
hors ligne pour LANCE}, 7 octobre 2026. \href{sans-llm-experiences.pdf}{Rapport}
et \href{sans-llm-experiences-sources.json}{registre des lectures}.
\bibitem{nist115} NIST, SP~800-115, \emph{Technical Guide to Information Security
Testing and Assessment}, septembre 2008, DOI 10.6028/NIST.SP.800-115.
\href{https://nvlpubs.nist.gov/nistpubs/Legacy/SP/nistspecialpublication800-115.pdf}{PDF officiel}.
Lecture ciblée des sections 3.3--3.4, pages imprimées 3-2--3-3 ; consultation
le 7 octobre 2026, sans lecture intégrale ni reproduction.
\bibitem{nist53a} NIST, SP~800-53A Rev.~5, \emph{Assessing Security and Privacy
Controls in Information Systems and Organizations}, janvier 2022. La notice
signale une mise à jour mineure 5.2.0 du 27 août 2025.
\href{https://csrc.nist.gov/pubs/sp/800/53/a/r5/final}{Notice officielle} et
\href{https://nvlpubs.nist.gov/nistpubs/SpecialPublications/NIST.SP.800-53Ar5.pdf}{PDF}.
Lecture du paragraphe 2.4, pages imprimées 11--12, et notice, le 7 octobre 2026 ;
pas d'affirmation de conformité ou de lecture intégrale de la révision.
\bibitem{lakens} D. Lakens, \emph{Equivalence Tests: A Practical Primer for
t Tests, Correlations, and Meta-Analyses}, Social Psychological and Personality
Science, 8(4), 355--362, 2017, DOI 10.1177/1948550617697177.
\href{https://pure.tue.nl/ws/portalfiles/portal/80918653/lakeequi2017.pdf}{PDF institutionnel}.
Lecture introduction et \emph{Testing for Equivalence}, pages imprimées 355--356,
le 7 octobre 2026. Transfert méthodologique proposé ; aucun test LANCE exécuté.
\end{thebibliography}
\end{document}
